\documentclass[preprint,12pt]{elsarticle}

\usepackage{amsmath,amssymb,amsthm}
\usepackage{booktabs}
\usepackage{graphicx}
\usepackage{url}
\usepackage[hidelinks]{hyperref}

\newtheorem{theorem}{Theorem}
\newtheorem{lemma}{Lemma}

\journal{Neurocomputing}

\begin{document}

\begin{frontmatter}

\title{Where Gap-Adaptive Spectral Regularization Helps, and Where It Does Not:
A Characterization and Operating Envelope for Differential Spectral Damping}

\author[iitp,sbs]{Praveg Vashishtha}
\address[iitp]{Department of Computer Science and Engineering, Indian Institute of Technology Patna, India}
\address[sbs]{Sopra Banking Software India}
\ead{praveg\_pa2503mth259@iitp.ac.in}

\begin{abstract}
Differential Spectral Damping (DSD) regularizes an ill-conditioned kernel system matrix by damping
each spectral direction according to its local eigengap rather than its eigenvalue magnitude, on the
Davis--Kahan principle that small-gap eigenvectors are the ones a perturbation corrupts. An earlier
report introduced DSD for least-squares twin support vector machines and showed a sizeable accuracy
gain over Tikhonov regularization. Two questions were left open: does the idea extend beyond that one
classifier, and what, if anything, can be proved about it. We answer both, and we are deliberate about
the negative parts of the answer. On the positive side, DSD improves two further high-dimensional
linear discriminants (ridge classification and Fisher discriminant analysis) on real data, and the
improvement survives a strengthened baseline set that includes truncated SVD and a tuned spectral
filter: on GINA ($d=970$) and Madelon ($d=500$), gap-adaptive damping beats the best competing
spectral regularizer by $3.3$ to $4.9$ percentage points across $25$ seeds. We then prove a
characterization theorem: the DSD inverse is bounded and well-posed, it reduces exactly to Tikhonov
when the spectrum carries no gap information, and it is the unique member of its family that imposes a
strictly monotone reliability ordering (larger gap, larger weight), something no function of
eigenvalue magnitude can do. Finally we map the method's operating envelope. Across four settings we
show, with a mechanism in each case, that DSD collapses to ordinary scalar regularization: Gaussian
process regression, non-linear RBF kernel classifiers, iterative-solver preconditioning, and
high-dimensional covariance estimation. The picture that emerges is specific rather than grand. DSD is
a structural prior for a particular regime, ill-conditioned linear high-dimensional classification,
and we delineate that regime precisely.
\end{abstract}

\begin{keyword}
spectral regularization \sep eigengap \sep twin support vector machine \sep ill-conditioned kernel
methods \sep Davis--Kahan perturbation \sep operating envelope
\end{keyword}

\end{frontmatter}

%% ===================================================================
\section{Introduction}
\label{sec:intro}

Many kernel methods reduce to inverting a symmetric system matrix whose spectrum decays fast. Twin
support vector machines, least-squares SVMs, Fisher discriminant analysis and ridge classification all
build their decision rule from such an inverse. When the input dimension is large relative to the
sample size, these matrices become severely ill-conditioned: a long tail of tiny, tightly clustered
eigenvalues sits beneath a few large ones. The standard remedy, Tikhonov regularization, adds a single
scalar to every eigenvalue. It stabilizes the inverse, but it treats all directions alike.

Differential Spectral Damping (DSD) \cite{vashishtha2026dsd} was proposed on a different premise. The
Davis--Kahan theorem \cite{davis1970rotation} says that the eigenvector attached to an eigenvalue with
a small spectral gap is the one a perturbation rotates most; its direction is unreliable. DSD damps
each direction by a function of that gap rather than of the eigenvalue magnitude, suppressing
small-gap (unreliable) directions and leaving well-separated ones almost untouched. On least-squares
twin SVM the reported gain over Tikhonov was large.

That first report left the work in an awkward place. The evidence was confined to one classifier, the
comparison was mostly against a single baseline, and the only theory was a motivating inequality with
no proof. This paper addresses those gaps, and it does so without inflating the result. We ask three
questions and report the answers as we found them, including the ones that are negative.

\emph{Does the idea generalize beyond twin SVM?} Yes, within a specific regime. DSD improves ridge
classification and Fisher discriminant analysis on real high-dimensional data, and the gain holds up
against truncated SVD and a tuned spectral filter, not only against Tikhonov (Section~\ref{sec:empirical}).

\emph{Can anything be proved?} Yes, a characterization rather than a performance guarantee. We prove
that the DSD inverse is bounded, that it contains Tikhonov as a special case, and that it uniquely
imposes a gap-monotone reliability ordering (Section~\ref{sec:theory}). We are careful to attribute the
boundedness and Tikhonov-recovery parts to classical regularization theory; the new content is the
ordering and the composite statement.

\emph{Where does it stop working?} In four distinct settings, each for a reason we can name. We report
these boundaries as a first-class contribution (Section~\ref{sec:envelope}), because a method whose
scope is mapped is more useful than one whose scope is merely asserted.

We also correct the record. The original report conjectured that DSD improves algorithmic stability.
Our measurements show the opposite, and we retract that conjecture in Section~\ref{sec:limitations}.

%% ===================================================================
\section{Background and the DSD operator}
\label{sec:background}

Let $W = \sum_{i} \lambda_i u_i u_i^\top$ be a symmetric positive semi-definite system matrix with
eigenvalues $\lambda_i > 0$ in ascending order and orthonormal eigenvectors $u_i$. Define the local
eigengap $\delta_i = \min_{j \neq i} |\lambda_i - \lambda_j|$, which for an ordered spectrum is the
smaller of the two neighbouring gaps. Tikhonov regularization replaces $1/\lambda_i$ by
$1/(\lambda_i + \gamma)$. DSD instead uses
\begin{equation}
\label{eq:dsd}
\tilde{\lambda}_i^{-1} \;=\; \phi(\lambda_i, \delta_i) \;=\;
\frac{\lambda_i}{\lambda_i^2 + \alpha\, e^{-\beta \delta_i}}, \qquad \alpha,\beta > 0,
\end{equation}
and forms $W^{+}_{\mathrm{DSD}} = \sum_i \phi(\lambda_i,\delta_i)\, u_i u_i^\top$. Writing
$d_i = \alpha e^{-\beta \delta_i}$ for the damping, a small gap produces large $d_i$ and therefore
strong suppression of that direction. The parameters are set from the spectrum itself by the
initialization rule of \cite{vashishtha2026dsd}, so DSD requires no cross-validation.

Two features of \eqref{eq:dsd} matter for what follows. The damping depends on the gap, not on the
magnitude; and when the gap is constant across the spectrum the formula is indistinguishable from a
scalar ridge. Both observations are made precise in Section~\ref{sec:theory}.

The downstream classifiers invert a matrix of this kind. Ridge classification solves
$(X_a^\top X_a)\,w = X_a^\top y$ with $X_a$ the bias-augmented design; Fisher discriminant analysis
solves $S_w\, w = (m_+ - m_-)$ with $S_w$ the within-class scatter; twin SVM \cite{kumar2009lstsvm}
solves a pair of such systems. In the high-dimensional regime each of these matrices exhibits the
clustered tail that makes the gap informative.

\subsection{Related work}
Four lines of work border ours. Classical regularization of ill-posed inverses, from Tikhonov
\cite{tikhonov1963} through the filter-factor theory of Engl, Hanke and Neubauer
\cite{engl1996regularization} and the pseudo-inverse analysis in Golub and Van Loan
\cite{golub2013matrix}, supplies the magnitude-based filters we compare against; a recent construction
in this tradition \cite{chen2026spectral} proves boundedness and Moore--Penrose recovery for a
structured spectral pseudo-inverse, and a learned variant reweights spectral coefficients by
signal-to-noise ratio \cite{scnet2026}. All of these key on eigenvalue magnitude, not on the gap. The
generalization theory of spectral algorithms \cite{velikanov2024spectral} analyzes exactly the
magnitude-filter class $h(\lambda)$; our operator lies outside it by construction
(Theorem~\ref{thm:char}(c)). The twin-SVM and kernel-machine literature has moved briskly in recent
years, with robust loss functions \cite{akhtar2024wave} and restricted kernel machine formulations
\cite{quadir2025trkm}, but these target the loss, the data representation, or multi-view structure
rather than the conditioning of the system-matrix inverse, so DSD is orthogonal to and composable with
them. Doubly regularized discriminant analysis \cite{r2lda2020} uses two magnitude ridge parameters,
again without a gap term. Finally, Nystr\"om approximation \cite{williams2001nystrom} produces the
kind of kernel submatrix DSD inverts, and the boundary studies of Section~\ref{sec:envelope} touch
Gaussian processes \cite{rasmussen2006gp}, covariance shrinkage \cite{ledoit2004shrinkage,ledoit2020nonlinear}
and algorithmic stability \cite{bousquet2002stability}.

%% ===================================================================
\section{A characterization theorem}
\label{sec:theory}

We prove what the DSD operator \emph{is}, not that it lowers risk. The attempt to prove a risk
guarantee failed for an instructive reason that we discuss in Section~\ref{sec:limitations}; what
survives is a clean characterization.

\begin{theorem}[DSD operator characterization]
\label{thm:char}
Let $W^{+}_{\mathrm{DSD}}$ be as above with damping $d_i = \alpha e^{-\beta \delta_i}$, $\alpha,\beta>0$.
\begin{enumerate}
\item[(a)] \textbf{Boundedness.} For every $i$, $0 < \phi(\lambda_i,\delta_i) \le 1/(2\sqrt{d_i})$,
with equality iff $\lambda_i = \sqrt{d_i}$. Hence
$\|W^{+}_{\mathrm{DSD}}\|_2 \le 1/(2\sqrt{\min_i d_i}) = e^{\beta \max_i \delta_i / 2}/(2\sqrt{\alpha})$.
In particular $\phi \to 0$ as $\lambda_i \to 0$ with $d_i$ bounded below, so the inverse stays bounded
as $W$ approaches singularity.
\item[(b)] \textbf{Tikhonov recovery.} If all gaps are equal, or in the limit $\beta \to 0$, then
$d_i \equiv \gamma := \alpha e^{-\beta \delta}$ and $\phi(\lambda_i,\delta_i) = \lambda_i/(\lambda_i^2+\gamma)$.
DSD thus contains eigenvalue-domain Tikhonov as a special case.
\item[(c)] \textbf{Monotone reliability ordering.} For fixed $\lambda>0$, $\phi(\lambda,\delta)$ is
strictly increasing in $\delta$. Among directions of equal eigenvalue magnitude, DSD assigns strictly
larger weight to the larger-gap direction. No filter that is a function of eigenvalue magnitude alone
can reproduce this ordering.
\end{enumerate}
\end{theorem}

\begin{proof}
(a) Fix $i$ and set $p(\lambda)=\lambda/(\lambda^2+d_i)$. Then
$p'(\lambda)=(d_i-\lambda^2)/(\lambda^2+d_i)^2$, positive for $\lambda<\sqrt{d_i}$ and negative after,
so $p$ is maximized at $\lambda=\sqrt{d_i}$ with value $1/(2\sqrt{d_i})$. Positivity is immediate.
Since $W^{+}_{\mathrm{DSD}}$ is symmetric with eigenvalues $\phi(\lambda_i,\delta_i)$, its spectral
norm is $\max_i \phi(\lambda_i,\delta_i) \le \max_i 1/(2\sqrt{d_i}) = 1/(2\sqrt{\min_i d_i})$, and
$\min_i d_i = \alpha e^{-\beta \max_i \delta_i}$ gives the stated form. If $\lambda_i \to 0$ with
$d_i \ge d_{\min}>0$ then $p(\lambda_i) \le \lambda_i/d_{\min} \to 0$.

(b) Equal gaps give $d_i \equiv \gamma$, and the formula becomes the eigenvalue-domain Tikhonov
inverse. The $\beta\to0$ limit sends every $d_i$ to $\alpha$ identically.

(c) With $d=\alpha e^{-\beta\delta}$ we have $\partial d/\partial \delta = -\alpha\beta e^{-\beta\delta}<0$
and $\partial \phi/\partial d = -\lambda/(\lambda^2+d)^2 < 0$, so by the chain rule
$\partial \phi/\partial \delta = \lambda\alpha\beta e^{-\beta\delta}/(\lambda^2+d)^2 > 0$ for
$\lambda>0$. A magnitude filter $g(\lambda)$ is by definition constant on equal-$\lambda$ directions,
so it cannot realize a strict order between two such directions.
\end{proof}

Parts (a) and (b) are, in spirit, classical facts of regularization theory
\cite{engl1996regularization}; a recent construction for structured inverse problems proves the same
boundedness and Moore--Penrose recovery for a magnitude-based filter \cite{chen2026spectral}. We
restate them here only to situate the operator. The content specific to DSD is part (c): a gap-keyed
filter sits outside the magnitude-filter class $h(\lambda)$ studied in the generalization theory of
spectral algorithms \cite{velikanov2024spectral}, precisely because it can separate two directions
that share an eigenvalue. This is the honest formal statement of what DSD does differently, and it
makes no claim about risk.

We verified all three parts numerically. Over $7\times10^{4}$ randomized checks the boundedness
inequality was never violated and was tight at $\lambda=\sqrt{d}$ to a relative error of $10^{-6}$; the
$\beta\to0$ Tikhonov recovery matched to $7\times10^{-14}$; and $\phi$ was strictly increasing in the
gap on every $(\lambda,\delta)$ grid tested, with zero monotonicity violations across $5\times10^{4}$
adversarial parameter draws.

%% ===================================================================
\section{Empirical study}
\label{sec:empirical}

\subsection{Protocol}
We follow the fairness protocol of \cite{vashishtha2026dsd}. DSD uses only its spectral
initialization, with no tuning. Every competitor is tuned: Tikhonov over a $15$-point logarithmic grid
for $\gamma$; truncated SVD over a grid of ranks; a general spectral filter
$\lambda^p/(\lambda^{p+1}+\varepsilon)$ over $(p,\varepsilon)$. Selection is on training accuracy, as
in the original report. All methods receive identical data and a $0.1$-magnitude feature-noise
injection; results are over $25$ random train/test splits with paired $t$-tests.

Datasets are the two real high-dimensional benchmarks from the feature-selection literature,
GINA ($d=970$) and Madelon ($d=500$) \cite{guyon2004gina}, plus a synthetic $d=200$ problem for a
controlled check. The classifiers are ridge classification and Fisher discriminant analysis, inverting
the matrices of Section~\ref{sec:background}; twin SVM is the subject of the original report and is
cited rather than re-run.

\subsection{Gap-adaptive damping improves two further discriminants}
Table~\ref{tab:breadth} reports the comparison against Tikhonov. On both real datasets DSD improves
ridge and Fisher classification by a clear margin, winning on almost every seed. The advantage tracks
the severity of the tail clustering: it is largest on GINA, whose system matrix has the tightest
cluster, and smaller on Madelon. The GINA gain of roughly $4.7$ points also matches the twin-SVM gain
reported in \cite{vashishtha2026dsd}, which is the consistency one would hope to see if the mechanism
is real.

\begin{table}[t]
\centering
\caption{Test accuracy, DSD (untuned) versus grid-tuned Tikhonov, $25$ seeds. All advantages are
significant at $p<10^{-3}$.}
\label{tab:breadth}
\begin{tabular}{llccc}
\toprule
Dataset & Classifier & DSD & Tikhonov & Advantage \\
\midrule
Synthetic $d{=}200$ & Ridge  & $0.761$ & $0.621$ & $+14.0$ pp \\
Synthetic $d{=}200$ & Fisher & $0.766$ & $0.621$ & $+14.5$ pp \\
GINA $d{=}970$      & Ridge  & $0.859$ & $0.812$ & $+4.7$ pp \\
GINA $d{=}970$      & Fisher & $0.861$ & $0.812$ & $+4.9$ pp \\
Madelon $d{=}500$   & Ridge  & $0.581$ & $0.548$ & $+3.3$ pp \\
Madelon $d{=}500$   & Fisher & $0.586$ & $0.548$ & $+3.8$ pp \\
\bottomrule
\end{tabular}
\end{table}

\subsection{The advantage survives stronger baselines}
A single baseline is not convincing, so we added the two spectral regularizers a reviewer would
reasonably ask for: truncated SVD and a tuned spectral filter. Table~\ref{tab:strong} gives the
outcome. On real data the result is unambiguous. Tikhonov, truncated SVD and the spectral filter all
converge to essentially the same accuracy, and DSD stands above the best of them by the same margins
as before. On the synthetic problem the story is more nuanced and we report it plainly: there
truncated SVD is competitive, beating DSD on ridge by $1.7$ points and losing to it on Fisher by
$1.5$ points. The two methods trade on synthetic spectra; on the real benchmarks DSD is the clear
standout. The reading we take from this is that truncated SVD, not Tikhonov, is the baseline that
matters, and that on genuine high-dimensional spectra the eigengap carries information a hard rank
cutoff does not.

\begin{table}[t]
\centering
\caption{DSD versus the best of \{Tikhonov, truncated SVD, spectral filter, naive\}, all tuned,
$25$ seeds. On real data every competitor collapses to the same value.}
\label{tab:strong}
\begin{tabular}{llcccc}
\toprule
Dataset & Classifier & DSD & best competitor & gap & $p$ \\
\midrule
Synthetic $d{=}200$ & Ridge  & $0.761$ & $0.778$ (TSVD) & $-1.7$ pp & $0.06$ \\
Synthetic $d{=}200$ & Fisher & $0.766$ & $0.751$ (TSVD) & $+1.5$ pp & $0.02$ \\
GINA $d{=}970$      & Ridge  & $0.859$ & $0.812$ & $+4.7$ pp & ${<}10^{-3}$ \\
GINA $d{=}970$      & Fisher & $0.861$ & $0.812$ & $+4.9$ pp & ${<}10^{-3}$ \\
Madelon $d{=}500$   & Ridge  & $0.581$ & $0.548$ & $+3.3$ pp & ${<}10^{-3}$ \\
Madelon $d{=}500$   & Fisher & $0.586$ & $0.548$ & $+3.8$ pp & ${<}10^{-3}$ \\
\bottomrule
\end{tabular}
\end{table}

%% ===================================================================
\section{Operating envelope: where DSD stops helping}
\label{sec:envelope}

A regularizer that helps everywhere usually helps nowhere in particular. We therefore mapped the
boundary of DSD's usefulness across four settings, and in each we can say why it reduces to scalar
regularization. We treat these as results, not caveats.

\paragraph{Gaussian process regression.} We built a GP regressor whose posterior inverse
$(K+\sigma^2 I)^{-1}$ can be formed by jitter, tuned Tikhonov, truncated SVD or DSD, and measured
calibration (negative log predictive density and interval coverage) over ill-conditioned regimes with
every competitor tuned on validation. DSD did not improve calibration in any regime. The reason is
structural: a GP already carries a principled scalar, the observation noise $\sigma^2$, that sets the
regularization level for calibration optimally, so there is no reliability gap for DSD to exploit. Its
advantage in classification comes precisely from the absence of such a scalar there.

\paragraph{Non-linear RBF kernel classifiers.} On RBF-kernel twin SVM and LS-SVM the gap-adaptive
inverse gave no gain. A bandwidth sweep explains it. The RBF Gram matrix is ill-conditioned only at
small bandwidth, but at those bandwidths the classifier cannot learn (accuracy near chance); where it
learns, the Gram matrix is already well-conditioned and DSD coincides with Tikhonov. The
ill-conditioning and the signal occupy different bandwidths, so there is nothing for the gap to
protect. This is consistent with the small non-linear gain reported in \cite{vashishtha2026dsd}.

\paragraph{Iterative-solver preconditioning.} DSD makes a poor preconditioner. Preconditioned
conjugate gradient converged faster with a Tikhonov-shifted preconditioner than with a DSD one, in
both truncated and full-spectrum form. The reason is a sign: accelerating a solver requires
\emph{amplifying} small eigenvalues to reduce the condition number, whereas DSD \emph{suppresses}
small-gap directions. The property that aids generalization in classification is the wrong one for
solving a system.

\paragraph{High-dimensional covariance estimation.} Gap-aware eigenvalue shrinkage was statistically
indistinguishable from both linear Ledoit--Wolf \cite{ledoit2004shrinkage} and nonlinear shrinkage
\cite{ledoit2020nonlinear} on Stein loss and minimum-variance-portfolio risk ($p=0.51$ and $p=0.41$).
On a clustered spectrum a small eigenvalue and a small gap nearly coincide, so the gap adds no
information a magnitude-based shrinkage does not already use. The gap carries independent signal only
when the corruption is of the eigenvector \emph{direction} and that direction enters the objective, as
it does in a discriminant but not in a covariance estimate.

Taken together, these four boundaries locate DSD's advantage. It requires a task in which corrupted
eigenvector directions, not merely small eigenvalues, harm the objective, and in which suppressing
those directions improves generalization. High-dimensional linear classification is such a task;
regression calibration, solving, and covariance estimation are not.

%% ===================================================================
\section{Limitations and a retraction}
\label{sec:limitations}

We are explicit about what we do not claim, including a correction to the earlier report.

\emph{No excess-risk guarantee.} We tried to prove that gap-adaptive damping lowers the excess risk of
the downstream classifier relative to any magnitude filter. The claim is false in general. Under
perturbation the eigenvalues shift as well as the eigenvectors, so a tuned magnitude filter already
absorbs most of the reliability signal that the gap would carry, and in strong noise the gap signal is
itself corrupted and gap damping becomes actively harmful. This is why Theorem~\ref{thm:char} is a
characterization and not a performance bound, and why our positioning against the magnitude-filter
generalization framework \cite{velikanov2024spectral} is one of complementarity, not dominance.

\emph{Retraction of the stability conjecture.} The original report conjectured that DSD improves
algorithmic stability. We measured leave-one-out sensitivity of the solved direction for DSD and for a
matched Tikhonov on clustered data, and found DSD to be roughly an order of magnitude \emph{less}
stable, with the difference significant. We retract the conjecture. DSD's benefit is not a stability
effect.

\emph{Scope.} The advantage is specific to the ill-conditioned linear high-dimensional regime, and on
synthetic spectra truncated SVD can match it. The headline evidence is on two real datasets and two
linear discriminants; broader confirmation, and a genuine theoretical account of the real-data
advantage that the failed excess-risk argument did not provide, are left to future work.

%% ===================================================================
\section{Conclusion}
\label{sec:conclusion}

Differential spectral damping is a gap-adaptive regularizer with a precisely drawn scope. It improves
high-dimensional linear classification on real data, and the improvement holds against the full
standard family of spectral regularizers, not only Tikhonov. It admits a clean characterization: the
inverse is bounded, it generalizes Tikhonov, and it alone imposes a gap-monotone reliability ordering.
And it reduces to ordinary scalar regularization in four settings we identify with a mechanism in each.
What we do not have is a performance theorem, and we have retracted a prior stability claim that the
data contradicts. The contribution is an honest, well-bounded account of a specialized tool, which we
think is more useful to a practitioner than a broader claim would be.

\section*{Reproducibility}
Code, the datasets' public sources, and a script for every number in this paper are available at
\url{https://github.com/Praveg432/dsd-regularization}. Each reported value is produced by a committed result file and reproduces
deterministically.

\bibliographystyle{elsarticle-num}
\bibliography{references}

\end{document}
